\section{The Poisson Process}
\label{subsection:Poisson_Process}

The Poisson process is the simplest continuous-time stochastic model of individual events. Imagine that people arrive at a front desk independently, with a constant arrival rate $\lambda>0$. Let $X(t)$ count the number of arrivals between time zero and time $t$, with $X(0)=0$. The possible events are 
\begin{equation}
 n\longrightarrow n+1 \quad\text{at rate }\lambda.
\end{equation}
For a sufficiently small time interval $dt$, we assume
\begin{align}
 \Pb\{X(t+dt)=n+1\mid X(t)=n\}&=\lambda\,dt+o(dt),\\
 \Pb\{X(t+dt)=n\mid X(t)=n\}&=1-\lambda\,dt+o(dt),
\end{align}
with probabilities of two or more arrivals in $dt$ of order $o(dt)$. The rate $\lambda$ has units of inverse time. The mean number of arrivals over an interval of length $t$ is $\lambda t$.

\subsection{From an ensemble of realizations to the master equation}
Imagine $M$ independent replicas of this counting process, all starting with zero arrivals. Let $N(n,t)$ be the number of replicas in which $n$ arrivals have occurred by time $t$. This is a \emph{statistical ensemble}. For a sufficiently large ensemble, the fractions of replicas undergoing each transition approach their corresponding probabilities. We can therefore establish a balance equation describing how the number of replicas in each category changes over a sufficiently short time interval. This equation becomes exact, to first order in \(dt\), in the infinite-ensemble limit. The expected balance over a short interval is
\begin{equation}
 N(n,t + dt) =\lambda\,dt\,N(n-1,t)+(1-\lambda\,dt)N(n,t)+o(dt)M,
\end{equation}
for $n\geq1$, with the convention $N(-1,t)=0$. The equation describes the
\emph{expected} numbers of replicas: for finite $M$, the actual category counts fluctuate. In the limit of a large independent ensemble, the fractions
$N(n,t)/M$ converge to probabilities $P(n,t)=\Pb\{X(t)=n\}$.
Taking the limit $dt\to0$ gives the \emph{forward master equation}
\begin{align}
 \frac{dP(n,t)}{dt}&=\lambda P(n-1,t)-\lambda P(n,t),&&n\geq1,
 \label{ME_Final_Poisson_n}\\
 \frac{dP(0,t)}{dt}&=-\lambda P(0,t),&&n=0.
 \label{ME_Final_Poisson}
\end{align}
These are coupled ordinary differential equations, one for each possible
number of arrivals. Their initial conditions are $P(0,0)=1$ and
$P(n,0)=0$ for $n\geq1$.

\subsection{Solving the master equation}
For zero arrivals,
\begin{equation}
 P(0,t)=e^{-\lambda t}.
\end{equation}
For one arrival, the equation is
\begin{equation}
 \frac{dP(1,t)}{dt}+\lambda P(1,t)=\lambda e^{-\lambda t}.
\end{equation}
Multiplication by the integrating factor $e^{\lambda t}$ yields
\begin{equation}
 \frac{d}{dt}\left[e^{\lambda t}P(1,t)\right]=\lambda,
 \qquad P(1,t)=\lambda t\,e^{-\lambda t}.
\end{equation}
Repeating the procedure (or proving it by induction) gives
\begin{equation}
 \boxed{P(n,t)=\frac{(\lambda t)^n}{n!}e^{-\lambda t},\qquad n\geq0.}
 \label{Eq_Poisson_Final}
\end{equation}
Thus, $X(t)$ has a Poisson \emph{distribution} with parameter $\lambda t$.
In particular,
\begin{equation}
 \E[X(t)]=\lambda t,\qquad \operatorname{Var}[X(t)]=\lambda t,
 \qquad\frac{d\E[X(t)]}{dt}=\lambda.
 \label{ODE_Average}
\end{equation}
The corresponding equation for the mean is exact for this linear process.

\subsection{An alternative derivation using a generating function}
Using the probability generating function introduced in
Section~\ref{sec:basic_pgf}, define
\begin{equation}
 G(z,t)=\sum_{n=0}^{\infty}P(n,t)z^n.
\end{equation}
Multiply each equation of the system given in Eqs (\ref{ME_Final_Poisson_n})-(\ref{ME_Final_Poisson}) by its corresponding factor $z^n$ and sum over all $n\geq0$:
\begin{equation}
 \frac{\partial G(z,t)}{\partial t}=\lambda(z-1)G(z,t).
 \label{GenFunc_Poisson}
\end{equation}
For each fixed $z$, this is an ordinary differential equation in time.
Its general solution is
\begin{equation}
 G(z,t)=G(z,0)e^{\lambda(z-1)t}.
\end{equation}
\emph{The initial condition, not normalization alone}, determines the
prefactor. Because $X(0)=0$ with probability 1, this is, $P(0,0)=1$, we have $G(z,0)=1$, so
\begin{equation}
 G(z,t)=e^{\lambda t(z-1)}
 =e^{-\lambda t}\sum_{n=0}^{\infty}\frac{(\lambda t)^n}{n!}z^n.
\end{equation}
Comparing coefficients of $z^n$ recovers Eq.~\eqref{Eq_Poisson_Final}.
Normalization, $G(1,t)=1$, provides a consistency check.
If instead the initial count has PGF $G(z,0)$, its distribution at time
$t$ is the convolution of the initial count distribution with the
Poisson$(\lambda t)$ distribution.

\subsection{Why this is a continuous-time Markov process}
The probability of the next arrival depends only on the rate $\lambda$,
not on how many arrivals occurred previously. Its waiting time $T$ satisfies
\begin{equation}
 \Pb\{T>t\}=e^{-\lambda t},\qquad \E[T]=\frac{1}{\lambda}.
\end{equation}
Hence the waiting time is exponential and memoryless. We will return to this
result when introducing exact stochastic simulation algorithms.

\subsection{Exercises}
\begin{enumerate}[label=\textbf{Exercise \arabic*.},leftmargin=*]
 \item \textbf{Forward equation.} Starting from the small-$dt$ transition
 probabilities, derive the master equation for $P(n,t)$ and write the
 boundary equation for $n=0$.
 \item \textbf{Normalization.} Show directly from the master equation that
 $d\sum_{n=0}^{\infty}P(n,t)/dt=0$.
 \item \textbf{Moments.} Multiply the master equation by $n$ and sum over
 all $n$ to derive $d\E[X(t)]/dt=\lambda$. Repeat for $n^2$ and find
 the variance.
 \item \textbf{Generating function.} Solve the equation for $G(z,t)$
 when $X(0)=m$ rather than zero, and interpret the result.
 \item \textbf{Two time intervals.} For a homogeneous Poisson process,
 find the probability that there are no arrivals in $(0,s]$ and exactly
 two arrivals in $(s,s+t]$. State the additional property of independent
 increments being used.
\end{enumerate}

